\chapter{Fájlok, feltöltés és médiatár}

\noindent\textbf{Tanulási út: Gyors út: opcionális ág.} Most olvasd, ha az alkalmazás uploadot fogad vagy alkalmazástulajdonú médiát szolgál ki.\par\medskip

\invariantblock{Biztonsági invariáns: a kliens fájlmetaadata nem megbízható}{A beküldött fájlnév és Content-Type leíró metaadat, nem bizonyíték biztonságos tárolásra vagy tartalomra. A storage maradjon AppFs/capability határ mögött, a kiszolgálási döntés pedig trusted detektálásból és policyból származzon.}

\diagnosticblockhu{Compiler diagnostic}{SEC-FILE-001}{Egy \code{Image} uploadból hiányzik a media-kiszolgálás előtti explicit \code{publish} transition.}{Használd a typed image upload formát \code{publish} jelöléssel; a raw/nem ellenőrzött upload maradjon private.}

A feltöltés request-, erőforrás-, filesystem-, authorization-, persistence- és serving-határokat érint. A Velran ezért külön kezeli az általános, AppFs-be streamelt \code{Upload} értéket és a publikálható médiához való szigorúbb \code{Image} típust.

\section{A mentális modell: három külön dolog}

Feltöltésnél tartsd külön a következőket:

\begin{enumerate}
  \item \textbf{kliens-metaadat}: eredeti fájlnév és a kliens által állított MIME-típus;
  \item \textbf{tárolás}: a runtime által választott, AppFs-en belüli szerveroldali fájl;
  \item \textbf{üzleti referencia}: amit az alkalmazás adatbázisban tárol és később felhasznál.
\end{enumerate}

A kliens fájlneve nem storage path, a kliens \code{Content-Type} fejléce pedig nem biztonsági bizonyíték. Ezt a két értéket legfeljebb megjelenítési vagy audit-metaadatként kezeld.

\section{Általános fájlfeltöltés: \code{Upload}}

A route deklarációja megmondja a runtime-nak, hogy multipart fájlt vár, és melyik statikus AppFs alkönyvtárba helyezze azt:

\begin{lstlisting}[language=Velran]
route uploadFile POST "/upload"
    upload file<Upload> to "uploads"
    public => uploadFile;
\end{lstlisting}

\newpage
Az action már nem a teljes fájlt kapja meg memóriában. Az \code{Upload} a feltöltés metaadatait adja át:

\begin{lstlisting}[language=Velran]
#[action]
fn uploadFile(ctx: ActionContext, file: Upload)
    -> Result<Redirect, PageError> {
    let savedPath = file.path;
    let originalName = file.filename;
    let contentType = file.contentType;
    let byteCount = file.bytes;
    return Ok(redirect(uploaded()));
}
\end{lstlisting}

A redirect typed \code{uploaded()} GET route helpert használ. Ennél fontosabb, hogy a raw \code{Upload} továbbra is private application data; attól még nem lesz public asset, hogy a handler látja a metadata-ját. Ez egyben fontos erőforrás-védelmi tulajdonság: a runtime streameli a multipart tartalmat, ezért egy nagyobb fájl nem válik automatikusan ugyanekkora request-szintű memóriaallokációvá.

\section{AppFs: az alkalmazás írható fájltere}

Az AppFs az alkalmazás számára kijelölt, korlátozott fájlrendszer-rész. Productionban külön data rootot adj neki; ne keverd a program binárisával, a statikus assetek könyvtárával vagy a rendszer más írható területeivel.

Példa szerverkonfiguráció:

\begin{lstlisting}
[storage]
data_root = "/srv/velran/data"
fs_mode = "rwc"
max_upload_bytes = 16777216
max_image_pixels = 40000000
\end{lstlisting}

A rövid capability-jelölések jelentése:

\begin{description}
  \item[\code{r}] olvasás;
  \item[\code{w}] írás;
  \item[\code{c}] létrehozás és az upload commitjához szükséges műveletek.
\end{description}

Feltöltéshez tipikusan \code{rwc} szükséges. A byte-limit operátori hard limit: az alkalmazáskód nem emelheti meg tetszőlegesen egy route kedvéért.

Linuxon az AppFs confinement célja, hogy az alkalmazás a kijelölt root alól path traversal, symlink vagy magic-link trükkel se tudjon kitörni. Ettől még a normál Unix ownership és permission beállítások is szükségesek: a service user csak azt a könyvtárat kapja meg írhatónak, amelyre valóban szüksége van.

\section{Mi történik egy feltöltés során?}

A biztonságos feltöltési útvonal lépései fejlesztői szemmel:

\begin{enumerate}
  \item a runtime ellenőrzi a request méretkorlátait és a route upload-sémáját;
  \item a fájlt streaming módon egy confined staging helyre írja;
  \item a storage kulcsot a szerver generálja, nem a kliens fájlnevéből képezi;
  \item typed feltöltésnél további szerveroldali tartalomellenőrzés fut;
  \item sikeres feldolgozás után a fájl atomikusan commitolható;
  \item handler- vagy runtime-hiba esetén a runtime cleanupot kísérel meg.
\end{enumerate}

Ne építs AppFs mellett párhuzamos fájlnév-szűrő írási útvonalat; az megkerülné a framework határát.

\section{Képfeltöltés: preferáld az \code{Image} típust}

CMS-ben, wikiben, hírportálon és adminfelületen a képekhez az általános \code{Upload} helyett az \code{Image} a megfelelő surface:

\begin{lstlisting}[language=Velran]
route saveHero POST "/admin/hero"
    upload hero<Image> to "media" publish
    auth role Editor
    => saveHero;
\end{lstlisting}

Az \code{Image} ugyanazt a streaming multipart és AppFs útvonalat használja, de a szerver a tényleges fájltartalmat is ellenőrzi. V1-ben PNG és JPEG fogadható el. A kliens által küldött MIME-típus nem authority.

Az SVG szándékosan nem általános képfeltöltési formátum: aktív tartalmat és script/XSS felületet hordozhat. A GIF, WebP vagy AVIF támogatását se feltételezd, ha az adott runtime-verzió dokumentációja nem mondja ki.

\section{Byte-limit és pixel-limit nem ugyanaz}

Képnél két külön erőforrás-probléma van. Egy fájl lehet kevés byte, miközben a dekódolt kép óriási pixelszámú. Ezért a \code{max_upload_bytes} mellett külön \code{max_image_pixels} limit is van.

\begin{lstlisting}
max_upload_bytes = 16777216
max_image_pixels = 40000000
\end{lstlisting}

A második limit a \code{width * height} értéket korlátozza. A production konfigurációban mindkettőt tudatosan állítsd be az alkalmazás feladatához; a ``16 MB-os upload limit'' önmagában nem kép-dekódolási védelem.

\section{Kép referencia az adatbázisban}

Az \code{Image} first-class scalar, ezért model mezőként és adatbázisban is tárolható canonical reprezentációval. Az alkalmazásnak nem kell -- és nem is érdemes -- a belső storage path részeit kézzel összeraknia vagy szétszednie.

Tipikus üzleti modell:

\begin{lstlisting}[language=Velran]
model Article {
    id: i64
    slug: Slug
    title: String
    body: String
    hero: Image
}
\end{lstlisting}

A lényeg az, hogy az adatbázisban \emph{referencia} legyen, ne kliens fájlnév és ne kézzel gyártott publikus URL. Így a tárolási és a HTTP-kiszolgálási szabályok a runtime kontrollja alatt maradnak.

\newpage
\section{Biztonságos megjelenítés HTML-ben}

Egy validált képet nem közvetlen string-interpolációval jelenítünk meg, hanem az erre szolgáló renderelővel:

\begin{lstlisting}[language=Velran]
#[page]
fn article(ctx: PageContext, article: Article)
    -> Result<Html, PageError> {
    return Ok(html {
        <article>
            <h1>{{ article.title }}</h1>
            @image(article.hero, article.title)
            @markdown(article.body)
        </article>
    });
}
\end{lstlisting}

Az \code{@image(image, alt)} a validált kép-referenciából készít HTML-t. Az alt szöveget escape-eli, a validált méreteket felhasználja, és a beépített read-only media endpointot célozza. Az \code{Image} közvetlen \code{{ image }} interpolációja compiler error, így a typed referencia nem csúszhat át véletlenül raw URL vagy markup szerepbe.

\section{A beépített media endpoint}

A runtime fenntartja a következő URL-prefixet:

\begin{lstlisting}
/__velran/media/
\end{lstlisting}

Az alkalmazás ezzel ütköző route-ot nem deklarálhat. A media endpoint nem teszi HTTP-n böngészhetővé a teljes \code{data_root} könyvtárat: csak olyan AppFs könyvtárból szolgál ki képet, amelyet a program valamely \code{upload ...<Image> to "..." publish} route-ja deklarált.

Kiszolgáláskor a szerver a tartalmat ismét a képdetektor alapján kezeli, nem a kliens eredeti MIME-jára hagyatkozik. A válaszhoz biztonsági és cache headerek társulnak; a média kiszolgálása nem futtat alkalmazás-route-ot és nem hoz létre sessiont.

\section{CMS-példa: hero kép feltöltése}

Egy szerkesztői felületnél a gyakorlati folyamat általában két üzleti lépésből áll:

\begin{enumerate}
  \item a szerkesztő feltölti a validált képet;
  \item az alkalmazás az így kapott typed képreferenciát az adott cikkhez rendeli.
\end{enumerate}

A route legalább autentikált legyen, admin/szerkesztő funkciónál pedig role policy is indokolt:

\begin{lstlisting}[language=Velran]
route saveHero POST "/admin/articles/:id<i64>/hero"
    upload hero<Image> to "media" publish
    auth role Editor
    => saveHero;
\end{lstlisting}

Ne abból indulj ki, hogy ``csak kép, tehát veszélytelen''. Az upload route ugyanúgy state-changing POST: auth, CSRF/origin védelem, request-limit és auditálható üzleti művelet tartozhat hozzá.

\section{Adatbázis és AppFs együtt alkothat üzleti állapotot}

Ha DB-sor AppFs-képre hivatkozik, a DB és AppFs együtt alkot helyreállítandó üzleti állapotot. Kontrollált stop/drain mellett készíts összerendelt application DB + local-auth DB + AppFs snapshotot; egymástól független mentések nem automatikusan konzisztensek. A Redis session/cache/rate-limit állapota nem üzleti source of truth.

\section{Statikus asset és feltöltött fájl nem ugyanaz}

A programhoz tartozó CSS, JavaScript, ikon és egyéb verziózott asset más életciklusú, mint a felhasználói feltöltés. A static root legyen read-only deploy-artifact; az AppFs data root pedig külön, kontrolláltan írható runtime adat.

Ez a szétválasztás több hibát megelőz:

\begin{itemize}
  \item deploy nem írja felül a felhasználói médiát;
  \item upload nem módosíthat alkalmazáskódot vagy verziózott assetet;
  \item backupnál egyértelmű, mi üzleti adat és mi újratelepíthető release-artifact;
  \item a reverse proxy szabályai is tisztábban kialakíthatók.
\end{itemize}

\section{Mit ne csinálj?}

\begin{itemize}
  \item Ne használd a kliens fájlnevét közvetlen filesystem pathként.
  \item Ne bízz a multipart \code{Content-Type} értékében.
  \item Ne tedd publikussá a teljes AppFs data rootot Apache/Nginx aliasszal.
  \item Ne interpolálj \code{Image} értéket kézzel URL-be vagy HTML-be.
  \item Ne kezeld a DB-t és az AppFs-t egymástól független recovery state-ként, ha vannak köztük referenciák.
  \item Ne kerüld meg az operátori upload- és pixel-limiteket alkalmazáskódból.
\end{itemize}

\section{A jelenlegi media surface határai}

A V1 media felülete tudatosan kicsi. Ne építs dokumentálatlan feltételezésekre. A jelenlegi alap nem ígér automatikus thumbnail/resize pipeline-t, WebP/AVIF konverziót, EXIF-normalizálást vagy metadata strippinget, teljes admin media browsert, S3/object-storage backendet vagy általános Markdown image syntaxot.

Ha az alkalmazásnak ezek valamelyike szükséges, azt külön architekturális feladatként kezeld, saját erőforrás-, security- és lifecycle-szabályokkal.

\section{Fejlesztői ellenőrzőlista feltöltéshez}

{\small
\begin{tabularx}{\textwidth}{@{}Y Y@{}}
1. Auth/authorization rendben? & 2. AppFs root elkülönített és nem publikus? \\
3. Ésszerű byte-limit? & 4. Képnél pixel-limit is? \\
5. Képhez \code{Image}? & 6. Filename/MIME csak metaadat? \\
7. DB-ben typed referencia? & 8. HTML-ben \code{@image}? \\
9. DB + AppFs együtt mentve? & 10. Nincs reverse-proxy data-root alias? \\
\end{tabularx}
}

\section{Ellenőrzött companion példák}
A pontos aktuális szintaxishoz használd:
\begin{itemize}
  \item \path{examples/upload/app.vrn} -- általános upload-határ;
  \item \path{examples/media/app.vrn} -- validált media felület;
  \item \path{examples/typed-multipart/main.vrn} -- típusos multipart dekódolás.
\end{itemize}

\section{Gyakorlat: threat-modellezz egy feltöltést}
\paragraph{Gyakorlási mód: támogatott.} Használd az előszóban meghatározott támogatott módszert.

Egy képfeltöltésnél sorold fel a byte-méretet, pixelszámot, fájlnevet, MIME metaadatot, dekódolt formátumot, tárolási helyet, autorizációt, backupot és serving útvonalat. Döntsd el, mely érték lesz megbízható csak framework-validáció után, és mely marad örökre nem megbízható metaadat.

\section{Tipikus hiba: uploadot teszünk a static root alá}
A static asset deployment artefaktum; az upload módosuló alkalmazásadat. A kettő összekeverése gyengíti a rollback és backup szemantikát, és könnyen megkerülhet autorizációs vagy tartalomkezelési szabályokat.

\section{Folyamatos mintaprojekt: Secure Notes}
Opcionálisan csatolj egy képet a jegyzethez a típusos media útvonalon. A fájl maradjon a static release-fán kívül, az adatbázis csak a szükséges típusos referenciát tárolja, és ellenőrizd, hogy egy jegyzet törlése vagy cseréje nem teheti elérhetővé más felhasználó médiáját.
